\documentclass{tufte-handout}

%\geometry{showframe}% for debugging purposes -- displays the margins



\newcommand{\bra}[1]{\left(#1\right)}
\usepackage{clrscode3e}
\usepackage{xcolor}
\usepackage{tikz}
\usepackage{amsmath,amsthm}
\usetikzlibrary{shapes}
\usetikzlibrary{positioning}
% Set up the images/graphics package
\usepackage{graphicx}
\setkeys{Gin}{width=\linewidth,totalheight=\textheight,keepaspectratio}
\graphicspath{{graphics/}}

\title{Math 455 notes}
\author{Shan Gao}
\date{Winter 2024}  % if the \date{} command is left out, the current date will be used

% The following package makes prettier tables.  We're all about the bling!
\usepackage{booktabs}

% The units package provides nice, non-stacked fractions and better spacing
% for units.
\usepackage{units}

% The fancyvrb package lets us customize the formatting of verbatim
% environments.  We use a slightly smaller font.
\usepackage{fancyvrb}
\fvset{fontsize=\normalsize}

% Small sections of multiple columns
\usepackage{multicol}

%--------Theorem Environments--------
%theoremstyle{plain} --- default
\newtheorem{theorem}{Theorem}
\newtheorem{corollary}[theorem]{Corollary}
\newtheorem{proposition}[theorem]{Proposition}
\newtheorem{lemma}[theorem]{Lemma}
\newtheorem{conjecture}[theorem]{Conjecture}
\newtheorem{question}[theorem]{Question}

\theoremstyle{definition}
\newtheorem{definition}[theorem]{Definition}
\newtheorem{definitions}[theorem]{Definitions}
\newtheorem{construction}[theorem]{Construction}
\newtheorem{example}[theorem]{Example}
\newtheorem{examples}[theorem]{Examples}
\newtheorem{notation}[theorem]{Notation}
\newtheorem{notations}[theorem]{Notations}
\newtheorem{addemdum}[theorem]{Addendum}
\newtheorem{exercise}[theorem]{Exercise}

\theoremstyle{remark}
\newtheorem{remark}[theorem]{Remark}
\newtheorem{remarks}[theorem]{Remarks}
\newtheorem{warning}[theorem]{Warning}
\newtheorem{scholium}[theorem]{Scholium}


% These commands are used to pretty-print LaTeX commands
\newcommand{\doccmd}[1]{\texttt{\textbackslash#1}}% command name -- adds backslash automatically
\newcommand{\docopt}[1]{\ensuremath{\langle}\textrm{\textit{#1}}\ensuremath{\rangle}}% optional command argument
\newcommand{\docarg}[1]{\textrm{\textit{#1}}}% (required) command argument
\newenvironment{docspec}{\begin{quote}\noindent}{\end{quote}}% command specification environment
\newcommand{\docenv}[1]{\textsf{#1}}% environment name
\newcommand{\docpkg}[1]{\texttt{#1}}% package name
\newcommand{\doccls}[1]{\texttt{#1}}% document class name
\newcommand{\docclsopt}[1]{\texttt{#1}}% document class option name

\begin{document}

\maketitle
\begin{abstract}
    Personal notes for math 455
\end{abstract}
\begin{remarks}
    We begin with some general remarks about Hilbert spaces. First of all if $\{e_i\}_{i=1}^\infty$ is orthonormal, then 
    $$\|f\|_2^2 = \sum_{i=1}^\infty \langle f, e_i \rangle^2$$
    Furthermore, TFAE: 
    \begin{itemize}
        \item $\{e_i\}_{i=1}^\infty$ forms a basis of $H$
        \item if $\langle f,e_i \rangle = 0$ $\forall i$, then $f = 0$
        \item Let $a_i = \langle f, e_i \rangle$ and $S_n = \sum_{i=1}^N a_i e_i$, then $\|f - S_n(f)\|_2^2 \to 0$ as $n \to \infty$
        \item Defining
        $$ f \to \sum_{i=1}^N \langle f, e_i \rangle_i = F_N(f)$$ as the projection of $f$ onto the space $\{e_1,\ldots,e_n\}$, then $\forall f \in H$, $\|f\|_2^2 = \sum_{i=1}^\infty |\langle f, e_i \rangle|^2 = \sum_{i=1}^\infty |a_i|^2$
    \end{itemize}
    Also finite linear combinations of $\{e_i\}_{i=1}^\infty$ are dense in $H$.
\begin{remark}
    A fact that will be useful for the homework is that given $f \in L^1(-\pi, \pi)$, then the $n$-th Fourier coefficient of $f$ is $\hat{f}(n) = \langle f(x),e^{inx} \rangle = \frac{1}{2\pi} \int_{-\pi}^\pi f(x)e^{-inx}dx$ where $n \in \mathbb{Z}$.
\end{remark}
\end{remarks}

\begin{theorem}[Hilbert space basis existence] \label{thm:1}
    Every Hilbert space has an orthornormal basis
\end{theorem}
The proof for this uses Zorn's lemma (axiom of choice)
\begin{lemma}[Zorns]
    Every non-empty partially ordered set in which every chain is bounded also contains a maximal element
\end{lemma}
\begin{definition}[Partially ordered set]
    $(X, \leq)$ is a partially ordered set if 
    \begin{itemize}
        \item $x \leq x$ $\forall x \in X$
        \item if $x,y \in X$ with $x \leq y$ and $y \leq x$, then $x = y$
        \item if $x,y,z \in X$, with $x \leq y$, $y \leq z$, then $x \leq z$
    \end{itemize}
    
\end{definition}

\begin{definition}[Chain]
    A chain $A$ is a totally ordered subset of $X$ where $\forall x,y \in A$, either $x \leq y$ or $y \leq x$. 
\end{definition}
\begin{definition}[Bounded chain]
   $Y \subseteq X$ is bounded if $\exists z \in X$ such that $\forall y \in Y$, $y \leq z$. 
\end{definition}

\begin{definition}[maximal element]
    A maximal element $m$ of $X$ is an element where if $m \leq x$ for $x \in X$, then $x = m$.
\end{definition}

We want to use Zorn's lemma to prove 
\begin{proof}[Proof of Theorem \ref{thm:1}]
    Let $H$ be a Hilbert space, and let $\mathcal{O}$ be the collection of all orthonormal subsets of $H$. $\mathcal{O}$ is non empty because $\{x\} \in \mathcal{O}$ where $x$ is just any element of $H$ with $\|x\| = 1$. We can define a partial ordering by set inclusion. Observing that every chain $\mathcal{A} \subseteq \mathcal{O}$ has an upper bound, by taking the union of every set in $\mathcal{A}$. Knowing this Zorn's lemma guarantees the existence of a maximal element, call it $B$, we want to show that this forms a basis of $H$. It's enough, as shown previously, to show that the linear span of $B$ is dense in $H$, and this is if and only if the closure of $\text{span}B$, denoted, $\overline{\text{span}B}$ ends up being the whole space. 
    
    \begin{remark}[Orthogonal decomposition theorem]
        Letting $H$ be a hilbert space and $M$ be a closed subset of $H$, define the orthogonal complement of $M$ as 
        $$M^{\perp} = \{ y \in X \colon \langle x,y \rangle = 0, x \in X\}$$ 
        then for any $x \in H$, $\exists! $ decomposition where 
        $$x = Px + Qx$$
        where $Px \in M$ and $Qx \in M^\perp$, where $P$ and $Q$ are the projection of $x$ to $M$ and $M^\perp$, respectively. Generally we have that 
        $H = M \bigoplus M^\perp$
        
    \end{remark}
    Thus we can take the closure of the linear span of $B$, $\overline{\text{span} B}$, by the orthogonal decomposition theorem, if $\overline{\text{span}B} \neq H$, then 
    $$ H = \overline{\text{span}B} + \overline{\text{span}B}^\perp$$
    But that means exists some $y \in \overline{\text{span}B}^\perp$ with $\|y\| = 1$ that is orthonormal to every element in $B$, since $B \subseteq \overline{\text{span}B}$, thus $B \cup \{y\}$ is a larger orthonormal subset, contradicting maximality of $B$, and we are done.
\end{proof}

\begin{definition}[Trigonometric Polynomials]
    Trigonometric polynomials are polynomials of the norm 
    $$  \sum_{n=-N}^{+N} a_n e^{inx} \quad (a_n \in \mathbb{C})$$

    
\end{definition}
\begin{proposition}
    Trigonometric polynomials are dense in $L^2([-\pi,\pi])$
    
\end{proposition}
Recalling that $C([-\pi,\pi])$ is dense in $L^2([-\pi,\pi])$. It suffices to show that for all $f \in C([-\pi,\pi])$, then $\forall \varepsilon >0$ $\exists P$ a trig polynomial such that $\|f - P\|_2 < \varepsilon$. In fact we can show that $\exists P$ such that $\|f - P\|_{\infty} < \varepsilon$, which is even a stronger result. 

\begin{theorem}[Stone-Waierstrass]
    Let $X$ be a compact, Hausdorff topological space. Let $\mathcal{A}$ be an algebra of continuous functions on $X$ that seperates points and contains constant functions. Then $\mathcal{A}$ is dense in $C(X)$.

    
\end{theorem}
Before any proofs, we define an algebra
\begin{definition}[Algebra (not set theoretic)]
    define $\mathcal{A}$ to be an algbera, which is a set with multiplication addition and scalar multiplication on a field $\mathbb{K}$ such that $f, q \in \mathcal{A}$, and $c_1, c_2 \in \mathbb{K} \implies c_1 f + c_2 g \in \mathcal{A}$, and $f \cdot g \in \mathcal{A}$. (Examples include polynomials, trigonometric polynomials $e^{inx}$)
    
\end{definition}
\begin{definition}[Seperates points]
    We say a collection $\mathcal{F}$ of functions on $X$ separates points if $\forall x,y \in X$, $x \neq y$, $\exists f \in \mathcal{F}$ such that $f(x) \neq f(y)$. 
    
\end{definition}
We ain't gonna give a full proof, but for now consider a more restricted result. 
\begin{theorem}
    $\forall f \in C([-\pi,\pi])$, $\forall \varepsilon >0$, $\exists P$ a trigonometric polynomial such that $\|f - P\|_\infty < \varepsilon$. 
\end{theorem}

\begin{proof}
    Let $Q_1(x), Q_2(x),\ldots$ be trigonometric polynomials such that 
    \begin{itemize}
        \item $Q_k \geq 0$ $\forall k \in \mathbb{N}$ $\forall x \in \mathbb{R}$
        \item $\frac{1}{2\pi} \int_{-\pi}^\pi Q_k(x) dx = 1$ $\forall k \in \mathbb{N}$
        \item Letting $\eta_k(\delta) = \sup_{\delta \leq |x| \leq \pi} Q_k(x)$. Then $\lim_{k \to \infty} \eta_k(\delta) = 0$
    \end{itemize}
    The higher level concept of the proof, is that given a sequence of trigonometric polynomials with these propreties, given any function $f \in C([-\pi.\pi])$, we can construct a function, that is also a trigonometric polynomial such that for all $y \in [-\pi,\pi]$, these two are aribitrarily close. 
    Notice that defining
    $$ P_k(x) = \frac{1}{2\pi} \int_{-\pi}^\pi f(x-y)Q_k(y) dy$$. Using a change of variables and using the periodicity of $e^{inx}$ we get that the latter integral is equal to  
    \begin{align*}
        P_k(x) = \frac{1}{2\pi} \int_{-\pi}^{ \pi} f(y) Q_k(x-y)dy &= \frac{1}{2\pi} \int_{-\pi}^\pi f(y) \sum_{n=-N}^{+N} a_{n,k} e^{in(x-y)}dy \\
        &= \sum_{n=-N}^{+N} a_{n,k} e^{inx} \frac{1}{2\pi} \int_{-\pi}^\pi f(y)e^{-iny}dy \\
        &= \sum_{n=-N}^{+N}a_{n,k} \hat{f}_n(y) e^{inx} 
    \end{align*}
    which it itself a trigonometric polynomial, and then we take any function $f \in C([-\pi,\pi])$. Knowing this, let $\varepsilon > 0$, since $f$ is continuous on a compact interval, then it's uniform continuous, meaning that $\forall x_1,x_2 \in [-\pi,\pi]$ if $|x_1-x_2| < \delta$, then $|f(x_1) - f(x_2)| < \varepsilon$  and then we just need to evaluate 
    $$|f(y) - P_k(y)| = \frac{1}{2\pi}\int_{-\pi}^\pi (f(y-x) - f(y))Q_k(y) dx = I_1 + I_2 $$, since $\int Q_k(x) = 1$, we can bring $f(y)$ into the integral, since it doesn't depend on $x$. 
    where 
    $$I_1 = \frac{1}{2\pi} \int_{-\delta}^\delta (f(y-x) - f(y))Q_k(y)dx$$
    and 
    $$I_2 = \frac{1}{2\pi} \int_{[-\pi,-\delta]\cup[\delta,\pi]} (f(y-x) - f(y))Q_k(y)dx$$
    Observing that 
    \begin{align*}
        I_1 &= \frac{1}{2\pi} \int_{-\delta}^\delta (f(y-x) - f(y))Q_k(y)dx \\
        &\leq \frac{1}{2\pi} \int_{-\delta}^\delta |f(y-x) - f(y)|Q_k(y)dx \\
        &\leq \frac{1}{2\pi} \int_{-\delta}^\delta \varepsilon Q_k(y)dx \leq \varepsilon 
    \end{align*}
    Since $|(y-x) - y| = |x| \leq \delta$ in the case of $I_1$
    Now we just need to take care of $I_2$. Note that $\exists M \in \mathbb{N}$ such that  $|f(y-x) - f(y)| \leq \sup_{a,b \in [-\pi,\pi]} |f(b) - f(a)| \leq 2M$ since $f$ reaches it's maxima on a compact interval. 
    \begin{align*}
        I_2 &= \frac{1}{2\pi} \int_{[-\pi,-\delta]\cup[\delta,\pi]} (f(y-x) - f(y))Q_k(y)dx \\
        &\leq \frac{1}{2\pi} \int_{[-\pi,-\delta]\cup[\delta,\pi]} 2M \eta_k(\delta)dx = 2M\eta_k(\delta)
    \end{align*}
    which vanishes to $0$ as $k \to \infty$, by the third proprety, thus proving that the difference vanishes, so we are done with the main part of the proof. 
    Now we are left to construct $Q_1(x), Q_2(x),\ldots$ with the propreties as described above. 
   We can just choose $ \displaystyle Q_k(x) = c_k \left(\frac{1 + \cos x}{2} \right)^k $, where $c_k$ is a constant such that the integral turns out to be $1$. 
   and also observe that, given $c_k$ chosen appropriately
   \begin{align*}
     1 &= \frac{c_k}{2\pi} \int_{-\pi}^\pi \left( \frac{1 + \cos x}{2}\right)^kdx \\
     &= \frac{c_k}{\pi} \int_{0}^\pi \left( \frac{1+ \cos x}{2}\right)^k dx \quad \text{ by symmetry} \\
     &> \frac{c_k}{\pi} \int_{0}^\pi \left( \frac{1 + \cos x}{2}^k \sin x\right) \\
     &= \frac{c_k}{n} \int_{1}^0 u^k (-2du)\\
     &=\frac{c_k}{n} \int_{0}^1 u^k (2du) \\
     &= \frac{2 c_k}{\pi} \cdot \frac{1}{k+1}\\
   \end{align*}
   Using the substitution $u = \frac{1+\cos x}{2}$ and $du = \frac{1}{2} \sin x dx$. 
   We then get that $c_k < \frac{\pi}{2} (k+1)$
   so $Q_k(x) < \frac{\pi}{2} (k+1) \left(\frac{1 + \cos x}{2}\right)^k \to 0$ for $x \neq 0$, so we recover proprety 3. 




\end{proof}
\begin{corollary}
    trigonometric polynomials are dense in $L^2([-\pi,\pi])$

   \end{corollary} 
    

\begin{proof}
    Trig polys are dense in $C([-\pi,\pi])$ under $\|.\|_\infty$. If $|P_k(x) - f(x)| < \varepsilon$, then we simply get that 
    \begin{align*}
        \|P_k - f\|_2 = (\frac{1}{2\pi} \int_{-\pi}^\pi |f(x) - P_k(x)|^2dx)^\frac{1}{2} < \sqrt{\varepsilon^2} = \varepsilon
    \end{align*}

\end{proof}

\begin{theorem}
	Let $X,Y$ be Banach spaces, $F \colon X \to Y$ bounded linear operator, and $F$ is surjective onto $Y$. Let $U$ be the open unit ball in $X$ and $V$ be the open unit ball. Then $\exists \delta > 0$ such that $F(U) \supset \delta V \implies F(B_r(x_0)) \supseteq B_{\delta r} (F(x_0))$
\end{theorem}
\begin{proof}
	$\forall y \in Y$, $\exists x \in X$, such that $F(x) = y$. If $\|x\| \leq k$, then we have something $y \in F(k\cdot U)$, $Y = \bigcup_{k \in \mathbb{N}} F(k \cdot U)$, $Y$ complete. By Baire Category theorem, $\exists k$ such that $\overline{F(k\cdot U)} \supseteq W$, for some $W$ non empty open set in $Y$. For every $y \in W$, $\exists \{x_i\}_i \subseteq k \cdot U$ such that $F(x_i) \to y$, this is by the fact that $y$ in the closure of $F(k\cdot U)$ and the fact that $F$ is onto. So let $k$ be sufficiently large such that $\overline{F(k \cdot U)}$ contains non empty interior, so open set $W$. Let $y_0 \in W$, $\exists \rho >0$ such that $B(y_0, \rho) \subseteq W$, then $y_0 + y \in W$, if $\|y\| < \rho$, then $\exists \{x_i^\prime\}_i, \{x^{\prime\prime}_i\}_i \subseteq k \cdot U$ such that $F(x_i^\prime) \to y_0$, $F(x_i^{\prime\prime}) \to y_0 + y$, $F(x_i^{\prime\prime}- x_i^\prime) \to y_0 + y - y_0 = y$.  
\end{proof}

Let $f \in L^1([-\Pi])$ where $\Pi = [-\pi,\pi]$. Recall Fourier coefficents of $f$ are given by 
$$
    \hat{f}(n) = \frac{1}{2\pi} \int_{-\pi}^\pi f(x)e^{-inx}dx
$$
We can also define a partial sum of the Fourier convolution
$$
    (S_n f)(x) = \sum_{k=-n}^n \hat{f}(k) e^{ikx}dx = \frac{1}{2\pi} \int_{-\pi}^\pi f(x) \underbrace{\left(\sum_{k=-n}^n e^{ikx}\right)}_{\text{Dirichlet Kernel $D_n(x)$}} dx 
$$


\bibliographystyle{plainnat}
\bibliography{notespt2}

\end{document}

